\documentclass[10pt,a4paper]{amsart}
\usepackage[margin=19mm]{geometry}
\usepackage{amsmath,amssymb,mathtools}
\usepackage[scr=rsfs,cal=euler]{mathalfa}
\usepackage{xfrac}
\usepackage[unicode,hidelinks,bookmarks=false]{hyperref}
\newcommand{\ie}{i.e.\ }
\newcommand{\cf}{cf.\ }
\newcommand{\resp}{respectively\ }
\newcommand{\Subsection}{\subsection}
\providecommand{\nobreakdash}{\nobreak\hbox{-}}
\setlength{\emergencystretch}{2em}
\multlinegap0pt
\usepackage{bm}
\usepackage{bbm}
\usepackage{dsfont}
\usepackage{xspace}
\usepackage{cancel}
\usepackage{mathtools}
\usepackage{amsthm}
\makeatletter
\@ifundefined{thm}{\newtheorem{thm}{Theorem}}{}
\makeatother
\newcommand{\Hcal}{\mathcal{H}} % define \Hcal here
\title[Continuously Frame-Convertible Sequences]{Continuously Frame-Convertible Sequences}
\author{Chad Berner}
\date{Source: arXiv:2605.02042v1 (CC BY 4.0)}
\begin{document}
\maketitle

\noindent\emph{Source and licence.} Continuously Frame-Convertible Sequences. Version arXiv:2605.02042v1 (CC BY 4.0), distributed under the Creative Commons Attribution 4.0 International licence, as recorded in the arXiv licence metadata for this exact version and shown on the abstract page https://arxiv.org/abs/2605.02042v1. Full rights notice: licenses/arxiv-2605.02042-CC-BY-4.0.txt.\par\smallskip

\noindent\textit{Editorial scope.} Counted unit: the Thm whose source label is \texttt{mainclassthm} in the article (source0.tex lines 139--158, source label \texttt{mainclassthm}), delivered together with the article's own complete original proof of it (source0.tex lines 159--189, one \texttt{proof} environment of 31 lines). No mathematics is written, repaired or re-derived: statement and proof are the article's own text line for line. Required context retained uncounted: none. Every definition, display and label of those lines is retained unchanged. Omitted from this package and not redistributed: 1--138, 190--560. Nothing inside a retained excerpt is omitted or altered: every retained body is delivered exactly as the article's lines are written, so no omission receipt of any kind accompanies this package. Citations: the retained excerpts carry no citation command at all, so no bibliography entry of the article is transcribed and no reference is redistributed. Reference adaptation: none was needed and none was made; every reference inside the delivered unit resolves inside it, so no unresolved reference or citation is produced. Printed slips kept literal: no printed word, symbol, hypothesis or number of the counted statement or of its proof was altered. Layout and class: the article's own class is \texttt{amsart} with packages including amsmath, amssymb, graphicx, biblatex; the wrapper is the lane's pinned \texttt{amsart} editorial wrapper at 10pt on A4, which restates the theorem-like environments the retained text needs and adds the article's own macro definitions that the retained text needs (\texttt{\textbackslash Hcal}), with extra wrapper packages (bm, bbm, dsfont, xspace, cancel, mathtools). The delivered numbering is the unit's own. Assets: no retained span contains a figure environment, a graphics-inclusion command, a PDF-page inclusion, a table or a TikZ picture; no image file is copied into this package, the package declares no asset dependency and no image file is redistributed with it. Licence: version arXiv:2605.02042v1 of this article is distributed under the Creative Commons Attribution 4.0 International licence, as recorded in the arXiv licence metadata for this exact version and shown on the abstract page https://arxiv.org/abs/2605.02042v1. A full rights notice is delivered at licenses/arxiv-2605.02042-CC-BY-4.0.txt. Characters: every retained excerpt is pure ASCII, so the delivered engine, XeLaTeX with its default Latin Modern text font, reports no missing character.
\begin{thm}\label{mainclassthm}
The following are equivalent:
\begin{enumerate}
    \item $\{f_n\}_{n=0}^{\infty}\subseteq \Hcal$ is a CFC sequence
    \item $\operatorname{Im}(A_{f_n})$ contains an infinite dimensional closed subspace.
    \item There is an infinite-dimensional subspace $D\subseteq \Hcal$ and $A>0$ such that
$$A\|f\|^2 \leq \sum_{n=0}^{\infty}|\langle f, f_n\rangle|^{2}<\infty $$
for all $f\in D$.
    
\end{enumerate}

Furthermore, assuming $(3)$,
\begin{enumerate}
    \item If $D$ is closed, $\{f_n\}$ is a SCFC sequence. Conversely, if $\{f_n\}$ is a SCFC sequence, $D$ may be chosen to be closed.
    \item If $\overline{D}=\Hcal$, $\{f_n\}$ is an ICFC sequence. Conversely, if $\{f_n\}$ is an ICFC sequence, $D$ may be chosen so that $\overline{D}=\mathcal{H}.$
    \item 
   $\left\{(TT^{*})^{\frac{1}{2}}f_n  \right\}\subseteq \overline{D}$ is a Parseval frame in $\overline{D}$ where $T\in B(C, \Hcal)$ and $C$ is a closed subspace of $\ell^{2}(\mathbb{N})$.
\end{enumerate}

\end{thm}

\begin{proof}
$(1)\implies (2)$

It follows that for all $f\in \Hcal$,
\begin{equation}\label{analysisopeq}
\frac{1}{\|B\|^{2}} \|B^{*}f\|^{2}\leq \|f \|^{2}=\sum_{n=0}^{\infty}|\langle B^{*}f, f_n\rangle|^{2}<\infty.  
\end{equation}

Therefore if $\{A_{f_n}(B^{*}g_k)\}$ is a Cauchy sequence, then $\{g_k\}$ is a Cauchy sequence. It follows by continuity of $B^{*}$ that $B^{*}g_k\to B^{*}g$ for some $g\in \Hcal$. As a result, it must be that $\{\langle B^{*}g_k, f_n\rangle\}_{n=0}^{\infty}\to \{\langle B^{*}g, f_n\rangle\}_{n=0}^{\infty}$ in $\ell^2(\mathbb{N})$ by uniqueness of limits. Now we have shown that $A_{f_n}|_{\operatorname{Im}(B^*)}$ is defined, injective, and has closed range. Furthermore, $\operatorname{Im}(A_{f_n}|_{\operatorname{Im}(B^*)})$ is infinite dimensional since it is clear that $\operatorname{Im}(B^{*})$ is infinite dimensional.

$(2)\implies (3)$

Let $C$ denote the infinite dimensional closed subspace contained in \newline $\operatorname{Im}(A_{f_n})$. Also define 
$$D:=\left\{f\in \overline{\operatorname{span}\{f_n\}}: A_{f_n}\in C \right\}.$$
Now note $A_{f_n}|_{D}$ is invertible onto $C$. To show that this operator is also closed, suppose $\{g_k\}\subseteq D$ and $\{A_{f_n}|_{D}(g_k)\}$ are Cauchy sequences. It follows since $A_{f_n}|_{D}$ is invertible and $C$ is closed that $A_{f_n}|_{D}(g_k)\to A_{f_n}|_{D}(f)$ for some $f\in D$. Now since $A_{f_n}|_{D}(g_k)\to A_{f_n}|_{D}(g)$ pointwise as well where $g$ is such that
$g_k\to g$, it follows that 
$f=g$ since $f,g\in \overline{\operatorname{span}\{f_n\}}$.
Therefore, since we have shown that $A_{f_n}|_{D}$ is a closed invertible operator with closed range, its inverse is continuous by the closed graph theorem and $(3)$ follows.

$(3)\implies (1)$

It follows that $(A_{f_n}|_{D})^{-1}$ whose codomain is extended to $\Hcal$, is bounded. We will denote its bounded extension by $A_{f_n}^{-1}$. Therefore, for all $f\in D,$
$$\|A_{f_n}A^{-1}_{f_n}\{\langle f, f_n\rangle\}\|^{2}=\|\{\langle f, f_n\rangle\}\|^{2}=\sum_{n=0}^{\infty}|\left\langle \{\langle f, f_n \right\rangle\},(A^{-1}_{f_n})^{*}f_n\rangle|^{2},$$ showing $\left\{(A^{-1}_{f_n})^{*}f_n \right\}\subseteq \overline{\operatorname{Im}(A_{f_n}|_{D})}$ is a Parseval frame in $\overline{\operatorname{Im}(A_{f_n}|_{D})}$. Therefore, $\left\{U(A^{-1}_{f_n})^{*}f_n \right\}\subseteq \Hcal$ is a Parseval frame in $\Hcal$ where $U: \overline{\operatorname{Im}(A_{f_n}|_{D})}\to \Hcal$ is any unitary.

Now assuming $(3)$, by the proof of $(3)\implies (1)$, if $D$ is closed, $(A^{-1}_{f_n})^{*}$ has closed range and therefore, $U(A^{-1}_{f_n})^{*}$ is surjective. Conversely if $\{Bf_n\}$ is a Parseval frame for some bounded and surjective $B$, then $\operatorname{Im}B^{*}$ is closed, which we can choose to be $D$ as in the proof of $(1)\implies (2)$.

Furthermore, if $\overline{D}=\Hcal$, 
by the proof of $(3)\implies (1)$, $(A_{f_n}^{-1})^{*}$ would be injective and hence $U(A^{-1}_{f_n})^{*}$ is injective. Conversely if $\{Bf_n\}$ is a Parseval frame for some bounded and injective $B$, then $\overline{\operatorname{Im}B^{*}}=\Hcal$, which we can choose to be $D$ as in the proof of $(1)\implies (2)$.

Finally, note that there is a unitary $U: \overline{Im(A_{f_n}|_{D})}\to \overline{D}$ where $$U(A_{f_n}^{-1})^{*}=[A_{f_n}^{-1}(A_{f_n}^{-1})^{*}]^{\frac{1}{2}}.$$
\end{proof}

\end{document}
